\subsection{A bounded rational Prym exponent gives a cyclic-cover criterion}
\label{sec:bounded-prym-exponent}

\paragraph{Finding A34 (the actual-exponent criterion).}
Let $C/k$ be a smooth projective geometrically connected genus-two
curve over a finite field of characteristic two, let $P\in C(k)$,
and let $\pi:D\to C$ be a geometrically connected cyclic etale cover
of odd degree $n>1$, with constant deck group
$G=\langle\sigma\rangle$ and a rational point $P'\in D(k)$ above $P$.
Put
\[
 B=\ker(\operatorname{Nm}:J_D\longrightarrow J_C)^0,
 \qquad M=\exp\bigl(B(k)[2^\infty]\bigr),
 \tag{BE.1}
\]
where the exponent of the trivial group is one. If
\[
                         n>(2M-1)^2,
 \tag{BE.2}
\]
then the covering formulation of A1 separates the point section
of $P$ from every distinct $Q\in C(k)$.
There is no rank-one hypothesis, no exclusion of a separable
elliptic double quotient, and no requirement that $n$ be prime.
The exponent in (BE.1) is measured over the stated field $k$;
it must be measured again after a constant extension.

The conclusion is a sufficient criterion. The final corollary
below realizes $M=2$ and $n=19$: the criterion applies to that
cover even though its relative class number is even. This is
a strict extension of the rational-unit test for a specified
pointed cover over its stated field; it makes no nonexistence
claim about other covers of the same pointed curve.

\paragraph{Recovered and published infrastructure.}
The field generated by deck translates of a small function, its
Castelnuovo--Severi genus bound, and the resulting free elliptic
action already occur in recovered successor A.53
\cite[A.53]{ASrun}. Its stated nonabelian contradiction does not
apply to the cyclic group here. We retain that infrastructure and
give the cyclic four-point and wild-ramification argument in full.
The published input for the field bound is Poonen,
\emph{Gonality of modular curves in characteristic $p$},
Proposition~2.1, Lemma~2.2 and Corollary~2.3, author p.~3
\cite{Poonen2007Gonality}.

\begin{proof}
\emph{1. An exact torsion order gives a separable function.}
Suppose that $P$ and a distinct $Q\in C(k)$ collide in the
covering formulation of A1. Applied to the pointed cover in the
statement, that formulation supplies $Q'\in D(k)$ above $Q$ with
\[
                    \delta=[Q'-P']\in J_D(k)[2^\infty].
\]
Only this direction of the covering implication is used, as in
A01 and recovered A1 B.1 \cite[B.1]{A1run}.

One has $\operatorname{im}(\sigma-1)=B$.
Indeed the image is connected and lies in the norm kernel.
Conversely, $[n]-\sum_{i=0}^{n-1}\sigma^i$ is divisible by
$\sigma-1$ in the integral group ring; on $B$ it is $[n]$,
which is surjective. Thus
$t=(\sigma-1)\delta\in B(k)[2^\infty]$.
Let $m$ be its exact order, with $m=1$ when $t=0$.
Then $m\mid M$ and $m$ is a power of two. There is a function
$f$ on $D_{\bar k}$ with
\[
 \operatorname{div}(f)
       =m\bigl(\sigma Q'+P'-Q'-\sigma P'\bigr),
                  \qquad d:=\deg(f)=2m.
 \tag{BE.3}
\]
The four points are distinct, because the deck action is free
and the two fibres of $\pi$ are disjoint. If $m>1$ and $f$ were
a square in $\bar k(D)$, half the divisor in (BE.3) would be
principal, contradicting the exact order of $t$. For $m=1$,
the simple poles already exclude a square. Over the perfect
field $\bar k$, this proves that $f$ is separable.

\emph{2. The conjugate functions generate a small intermediate field.}
Set
\[
 K=\bar k(D),\qquad F=\bar k(\gamma f:\gamma\in G),
       \qquad e=[K:F],
\]
and let $Z$ be the smooth projective curve with function field $F$.
All generators have degree $d$. The cited Castelnuovo--Severi
consequence gives
\[
                 e\mid d,\qquad g(Z)\le(d/e-1)^2.
 \tag{BE.4}
\]
Precisely, Poonen's Corollary~2.3 applies because $F$ is
generated by degree-$d$ functions; its controlled-field
integer is $d/e$. Every extension here lies inside the
separable extension $K/\bar k(f)$.

Since $D\to C$ is etale, $g(D)=n+1$.
By (BE.2), $n>(d-1)^2$, so (BE.4) excludes $e=1$.
Thus $e>1$ is a power of two. The map $v:D_{\bar k}\to Z$
is separable of degree $e$, and $f=h\circ v$ for a nonconstant
function $h$ on $Z$.

\emph{3. The induced deck action on the intermediate curve is free.}
The field $F$ is $G$-stable. A kernel of the induced action
would have order dividing both $n$ and $e$, so the action on
$Z$ is faithful. It is free as well. A nontrivial stabilizer
would contain a subgroup of odd prime order $\ell$. That
subgroup acts freely on $D$ and preserves the fibre divisor
of $v$ above the fixed point. Its points are partitioned into
orbits of length $\ell$, with multiplicities constant on each
orbit. This would imply $\ell\mid e$, a contradiction.

Riemann--Hurwitz for $Z\to Z/G$ now gives
\[
                       g(Z)-1=n\bigl(g(Z/G)-1\bigr).
\]
The bound $0\le g(Z)<n$ from (BE.4) forces
$g(Z)=g(Z/G)=1$.
For the separable map $v$ its different divisor therefore has
degree
\[
                    \deg\operatorname{Diff}(v)=2g(D)-2=2n.
 \tag{BE.5}
\]

\emph{4. Wild ramification forces $e=2$.}
Choose a zero $z$ of $h$, and write $s=\operatorname{ord}_z(h)>0$.
Every point in the fibre $v^{-1}(z)$ must be one of the two
positive points in (BE.3), and the multiplicity of $f$ at either
one is $m$. Thus the fibre has one or two distinct points,
each with the same ramification index $m/s$, a power of two.
If $e>2$, that index is at least two. It is wild, and its
different exponent is at least its ramification index.

Here the last elementary local fact can be seen directly.
For uniformizers $u$ downstairs and $t$ upstairs over the
algebraically closed constant field, write
$u=a_qt^q+a_{q+1}t^{q+1}+\cdots$, where $q$ is the local
ramification index. If $q$ is even, the derivative of the
leading term vanishes, so
$\operatorname{ord}_t(du/dt)\ge q$; this derivative is nonzero
by separability. Its valuation is the different exponent.

Consequently the different contribution over $z$ is at least
the sum of the ramification indices, namely $e$.
The $n$ translates of $z$ under $G$ are distinct, since the
action on $Z$ is free. Equivariance gives the same lower bound
over each of them. If $e>2$, this would give
$\deg\operatorname{Diff}(v)\ge ne>2n$, contradicting (BE.5).
Hence $e=2$.

\emph{5. The two positive and two negative points exclude $e=2$.}
If the positive points $P'$ and $\sigma Q'$ had different
images in $Z$, the fibre over either image could contain
no other point: every such point would be an additional
zero of $f$. Since $v$ has degree two, both points would
be ramified. Their $G$-orbits are disjoint and each has $n$
points, because they lie above distinct points of $C$.
Every point in these two orbits has different exponent at
least two. This would give different degree at least $4n$,
again contradicting (BE.5).

Thus $P'$ and $\sigma Q'$ are the two distinct unramified
points of a single fibre. Ramification is invariant under
the deck action, so $Q'$ and $\sigma P'$ are unramified too.
Their pole fibres must likewise coincide: an unramified
degree-two fibre cannot be supported on only one pole.
Putting $x=v(P')$ and $y=v(Q')$, equivariance gives
\[
                         x=\sigma y,\qquad y=\sigma x.
 \tag{BE.6}
\]
Hence $\sigma^2x=x$. Since $n>1$ is odd, $\sigma^2$ is a
nonidentity element of $G$, contrary to the free action on
$Z$. This final contradiction proves A34.
\end{proof}

\paragraph{An exact way to measure the exponent at an inert cyclotomic prime.}
There is a useful arithmetic specialization, which does not
produce a cover. Let now $C/\mathbf F_Q$ have genus two in
characteristic $p$, let $P$ be rational, and let a constant
$P$-split cyclic cover have odd prime degree $r\ne p$, with
$\operatorname{ord}_r(p)=r-1$. Write $b:J_C(\mathbf F_Q)\to
\mathbf Z/r$ for its rational quotient and let $\chi$ be a
faithful character. The genus-two Riemann--Roch calculation
of A29 gives
\[
 L_\chi(T)=1+c_\chi T+Q\chi(b(\kappa_P))T^2,
 \quad c_\chi=\sum_{R\in C(\mathbf F_Q)}\chi(b([R-P])),
 \quad \kappa_P=[K_C-2P].
 \tag{BE.7}
\]
Put $E=\mathbb Q(\zeta_r)$ and normalize the valuation at
its unique prime above $p$ by $v(p)=1$. Then the rational
$p$-primary group of the whole Prym has the exact form
\[
 B(\mathbf F_Q)[p^\infty]\simeq
       \mathbb Z_p[\zeta_r]/p^e\mathbb Z_p[\zeta_r],
       \qquad e=v\bigl(L_\chi(1)\bigr),
 \tag{BE.8}
\]
including the zero group when $e=0$. Its exponent is $p^e$.

Here is a proof that avoids assuming a rank-one crystalline
comparison. The primitive-root condition makes
$\mathcal O=\mathbb Z_p[\zeta_r]$ the ring of integers of
the unramified extension of degree $r-1$. It acts on $B$,
because $1+\sigma+\cdots+\sigma^{r-1}$ vanishes on
$B=\operatorname{im}(\sigma-1)$. A finite $\mathcal O$-module
is a direct sum of modules $\mathcal O/p^{a_i}$.
Each nonzero summand contributes $r-1$ dimensions to its
subgroup killed by $p$. But
$B(\mathbf F_Q)[p]\subset B(\overline{\mathbf F}_p)[p]$,
whose dimension is the geometric $p$-rank, at most
$\dim B=r-1$. Thus at most one summand occurs.
Equivalently, the geometric etale Tate module is free over
$\mathcal O$ of rank zero or one; the finite-module argument
already suffices for the exponent assertion.

Finally the character product and the Jacobian norm isogeny
give
\[
 \#B(\mathbf F_Q)=N_{E/\mathbb Q}\bigl(L_\chi(1)\bigr).
\]
Since $p$ is inert and unramified in $E$, its valuation is
$(r-1)v(L_\chi(1))$. Comparing this with the order
$p^{(r-1)e}$ of the unique possible module proves (BE.8).
The same reasoning also proves that the exponent is one
when the rational $p$-primary group is zero.

\paragraph{The Weil valuation bound and its exceptions.}
Every conjugate of $c_\chi$ has absolute value at most
$2\sqrt Q$. Indeed its degree-two character factor is a
factor of the zeta numerator of the smooth etale cover,
whose Frobenius eigenvalues have absolute value $\sqrt Q$;
see \cite[Theorem~1.6]{WeilI}. Hence, putting $M_p=p^e$,
the algebraic-integer norm argument always gives
\[
                         M_p\le(\sqrt Q+1)^2.
 \tag{BE.9}
\]
For characteristic two and $Q\ge8$ it gives the sharper
dichotomy
\[
 \begin{cases}
 c_\chi=-1,\quad b(\kappa_P)=0,
                  \quad L_\chi(1)=Q,\quad M_2=Q;\\[2pt]
 \text{or}\qquad M_2\le2\sqrt Q+1.
 \end{cases}
 \tag{BE.10}
\]
The exceptional first alternative has all faithful
$L$-polynomials equal to $1-T+QT^2$.

For completeness, write $Q=p^a$.
Since $p$ is the unique prime over $p$, divisibility by
$p^e$ in the local ring means divisibility by the rational
integer $p^e$ in the ring of integers of $E$.
If a nonzero algebraic integer has all conjugates of
absolute value less than one, its nonzero integral norm
would have absolute value less than one, which is impossible.
Applying this to $L_\chi(1)/p^e$ proves (BE.9).
If $e\le a$, then $p^e$ divides $1+c_\chi$. Unless
$c_\chi=-1$, the same argument gives $p^e\le2\sqrt Q+1$.
If $e>a$ and $Q>2\sqrt Q+1$, applying it instead to
$(1+c_\chi)/Q$ forces $c_\chi=-1$; then
$L_\chi(1)=Q\chi(b(\kappa_P))$ has valuation $a$, a
contradiction. For $p=2$, $Q\ge8$ ensures the required
strict inequality. When $c_\chi=-1$, all its conjugates,
including $c_{\chi^{-1}}$, equal $-1$. The exact identity
$c_\chi=\chi(b(\kappa_P))c_{\chi^{-1}}$ from A29 forces
$b(\kappa_P)=0$, proving the exceptional alternative.
For $Q=2,4$, (BE.9) only gives the respective power-of-two
bounds $M_2\le4,8$; these small-field exceptions must not
be silently absorbed into (BE.10).

The exceptional alternative is already restricted by the
split fibre. Let $N=\#C(\mathbf F_Q)$ and let $N_i$ be its
Abel-image occupancies in the chosen quotient $\mathbf Z/r$.
The equality $1+\sum_{i=0}^{r-1}N_i\zeta_r^i=0$ implies
\[
 1+N_0=N_1=\cdots=N_{r-1}=\lambda,\qquad
 N+1=\lambda r.
 \tag{BE.11}
\]
Since $P$ occupies the zero coset, $N_0\ge1$ and
$\lambda\ge2$, so $r\le(N+1)/2$. If $N$ is even, then
$\lambda$ is odd and $\lambda\ge3$, giving
$r\le(N+1)/3$.

\paragraph{What the uniform valuation estimate does not improve.}
For $Q\ge8$ outside the exceptional case, replacing the
actual exponent in A34 by the bound in (BE.10) gives the
sufficient but coarse inequality
\[
                         r>(4\sqrt Q+1)^2.
 \tag{BE.12}
\]
This yields no improvement over the earlier B.3 size test.
To check that claim exactly, write $J=\#J_C(\mathbf F_Q)$,
$N=\#C(\mathbf F_Q)\ge1$, and let $a_2$ be the middle
coefficient of its genus-two zeta numerator. Weil gives
$N\le Q+1+4\sqrt Q$ and $a_2\le6Q$, while direct evaluation
of the numerator gives
\[
             J=(Q+1)N+a_2-2Q
                   \le(Q+1)N+4Q\le(5Q+1)N.
 \tag{BE.13}
\]
Condition (BE.12) implies both $r>N$ and $r>10Q+2$.
Consequently $N>2J/r$. If the rational quotient exists,
its kernel has order $J/r$, so both hypotheses
$r>N$ and $N\ge2|\ker b|$ of recovered A1 B.3 hold
\cite[B.3]{A1run}. In the even-$N$ case A26 already uses
the simpler inequality $r>N$.

For $Q=2,4$, the coarse bounds $M_2\le4,8$ would demand
$r>49,225$, respectively, whereas the Weil bound gives
$J<34$ and $J\le81$, respectively; no such rational
quotient can exist. In the exceptional case $M_2=Q$,
(BE.11) and the Weil bound give
$r\le(Q+2+4\sqrt Q)/2<(2Q-1)^2$ for $Q\ge2$, so A34's
inequality cannot hold there either.

Thus the potential gain is in measuring the actual rational
exponent in (BE.1), not in inserting only a uniform archimedean
bound. A positive density of primitive-root primes does not
give a constant pointed quotient at one of those primes, nor
does it supply a bound for this exponent after the cover and
its decks have been defined. Passing to a constant-deck field
$\mathbf F_{Q^d}$ replaces the square-root term by $Q^{d/2}$.
Before that passage Frobenius is semilinear for the deck action,
so the constant cyclic quotient and character-polynomial
argument above cannot simply be used over the original field.
No uniform control of that field degree or actual exponent is
proved here; no proof through Hooley follows from A34 alone.

\paragraph{A realized exponent-two witness of degree nineteen.}
The rank-one curve of A27 gives a strict example for the
actual-exponent criterion. Write
\[
 k_0=\mathbf F_{64}=\mathbf F_2[z]/(z^6+z+1),\qquad
 u=z^5+z^4+z^3+z,
 \qquad u^2+u+1=0.
\]
In the reciprocal coordinates $X=1/t$, $Y=X^2y$, its equation
and its rational Weierstrass base point are
\[
 C:\quad Y^2+X^2Y=uX^5+X^3+uX,
              \qquad W=\infty.
 \tag{BE.14}
\]
Thus $W$ is the simple-pole point of the original A27 model.
The two elliptic Frobenius factors of A27, with traces $1,-2$
over $\mathbf F_4$, have third power sums $-11,16$.
Consequently
\[
 \#C(k_0)=60,\qquad
 \#J_C(k_0)=(65+11)(65-16)=3724=19\cdot196.
 \tag{BE.15}
\]
All field elements in the following certificate are encoded
by their binary coefficients in the displayed polynomial basis.
The reduced Mumford divisor
\[
 \Delta=\bigl(X^2+54X+8,\;49X+50\bigr)
 \tag{BE.16}
\]
has exact order nineteen. Multiplication by $196$ defines
the surjective quotient $b:J_C(k_0)\to\mathbf Z/19$ by
$[196]A=b(A)\Delta$. The occupancies of the rational points
under $R\mapsto b([R-W])$, in residue order $0,\ldots,18$, are
\[
 (6,2,2,2,5,2,5,2,2,5,5,2,2,5,2,5,2,2,2).
 \tag{BE.17}
\]
These finite statements are checked by two separately
implemented divisor calculations: binary-linear interpolation
and extended-GCD Cantor composition. Both retain every point
image, the entire nineteen-element subgroup, and the exact
finite-field and cyclotomic arithmetic. The second calculation
uses cyclic length-nineteen convolution, independently of the
first calculation's polynomial-remainder multiplication.

Pull back the Lang quotient belonging to $\ker b$ along the
Abel map $R\mapsto[R-W]$. The connectedness and split-fibre
argument of A12 gives a geometrically connected constant
cyclic cover $D\to C$ of degree nineteen, with a rational
point above $W$. Let $B$ be its whole Prym, of dimension
eighteen. Since $W$ is Weierstrass, $\kappa_W=0$, and a
faithful character has
\[
 L_\chi(T)=1+cT+64T^2,
      \qquad c=\sum_{j=0}^{18}N_j\zeta_{19}^j.
 \tag{BE.18}
\]
The same two calculations give
\[
 \operatorname{ord}(\bar c)=7,
 \qquad 1+c^7\equiv2\pmod4,
 \qquad 1+c^7\equiv6\pmod8
 \quad\text{in }\mathbb Z[\zeta_{19}].
 \tag{BE.19}
\]
Here $\bar c$ is reduced modulo two. The order calculation
uses $\operatorname{ord}_{19}(2)=18$, so $\Phi_{19}$ is
irreducible modulo two; in particular the local coefficient
ring $\mathcal O=\mathbb Z_2[\zeta_{19}]$ is unramified of
degree eighteen. Only the congruence modulo four is needed
for the exponent conclusion.

Set $k=\mathbf F_{64^7}$. To justify the extension step
without assuming a unit-root comparison, let
$L_\chi(T)=(1-\alpha T)(1-\beta T)$ and put
$s_j=\alpha^j+\beta^j$. Then
\[
 s_0=2,\quad s_1=-c,\quad
 s_j=-c\,s_{j-1}-64s_{j-2}.
 \tag{BE.20}
\]
Modulo four this gives $s_7\equiv(-c)^7$.
The character factor over $k$ is
$1-s_7T+64^7T^2$, so (BE.19) gives
\[
 L_{\chi,k}(1)\equiv1+c^7\equiv2\pmod4.
\]
Formula (BE.8), applied over this stated field, now proves
\[
 B(k)[2^\infty]\simeq\mathcal O/2\mathcal O
                  \simeq(\mathbf Z/2)^{18},
 \qquad M=2,
 \qquad v_2\bigl(\#B(k)\bigr)=18.
 \tag{BE.21}
\]
Thus the relative class number is even, whereas
$19>(2\cdot2-1)^2=9$. A34 separates the section of $W$
from every distinct rational point of $C_k$ using this same
cover. The rational-unit criterion rejects this particular
cover over $k$, so the asserted strict extension is realized.
This does not assert that every other cover over $k$ fails
the unit criterion, or that a uniform supply of such
bounded-exponent covers has been proved.
